
\subsubsection{Dolne ograniczenia liczby gordyjskiej}
Dokładna wartość liczby gordyjskiej jest znana tylko dla niektórych klas węzłów, na przykład torusowych (fakt \ref{prp:torus_unknotting_number}) albo skręconych.
\index{węzeł!torusowy}%
\index{węzeł!skręcony}%

Borodzik oraz Friedl podali niedawno całkiem mocne ograniczenia na liczbę gordyjską w~pracach \cite{borodzik2014} i~\cite{borodzik2015}.
\index[persons]{Borodzik, Maciej}%
\index[persons]{Friedl, Stefan}%
Ich narzędziem jest parowanie Blanchfielda.
\index{parowanie Blanchfielda}%
Poprawiają tam starsze estymaty wynikające z~sygnatury Levine'a-Tristrama, indeksu Nakanishiego oraz przeszkody Lickorisha.
\index{indeks!Nakanishiego}%
\index{przeszkoda Lickorisha}%
\index{sygnatura!Levine'a-Tristrama}%
% DICTIONARY;Lickorish obstruction;przeszkoda Lickorisha;-
Wśród pierwszych węzłów o~co najwyżej 12 skrzyżowaniach dwadzieścia pięć ma liczbę gordyjską równą co najmniej trzy, trudno było uzasadnić to innymi metodami.

Dla satelitów znamy tylko bardzo słabe ograniczenia.
Z~pracy Scharlemanna i~Thompsona \cite{thompson1988} wynika, że jeśli wzorzec $P$ jest nietrywialny, a~jego homologiczna liczba owinięć (wartość bezwzględna klasy $[P] \in H_1(S^1 \times D^2) \cong \Z$) jest niezerowa, to satelita nietrywialnego węzła $K$ spełnia $\unknotting(K_P) \ge 2$.
\index{węzeł!satelitarny}%
\index[persons]{Scharlemann, Martin}%
\index[persons]{Thompson, Abigail}%
Hom, Lidman, Park \cite{hom2022} zaproponowali wzmocnienie tego wyniku:

\begin{conjecture}
    Niech $P$ będzie nietrywialnym wzorcem o~homologicznej liczbie owinięć $w \neq 0$, zaś $K$ nietrywialnym węzłem.
    Wtedy $\unknotting(K_P) \ge w + 1$.
\end{conjecture}

Nie można oczekiwać, że prawa strona będzie równa $w^2$: pewien kabel trójlistnika o~liczbie owinięć $2$ ma liczbę gordyjską co najwyżej $3 < 2^2$ \cite[problem 1.4]{baykur2026}.
\index[persons]{Hom, Jennifer}%
\index[persons]{Lidman, Tye}%
\index[persons]{Park, Junghwan}%
Ci sami autorzy twierdzą, że dla dowolnego nietrywialnego $K$ każdy kabel o~liczbie owinięć $w$ ma liczbę gordyjską co najmniej $w$.

Między genusem i~liczbą gordyjską nie ma ogólnej zależności: znamy węzły, dla których genus jest większy od liczby gordyjskiej, i~takie, dla których jest mniejszy.
W~niektórych klasach węzłów można jednak spodziewać się nierówności \cite[problem 1.5]{baykur2026}.
Węzeł jest dodatni, jeśli ma diagram, w~którym wszystkie skrzyżowania są dodatnie.
Z~prac Rudolpha \cite{rudolph1993, rudolph1999} wynika, że dla węzłów dodatnich zachodzi $\genus K \le \unknotting K$.
\index[persons]{Rudolph, Lee}%
\index{węzeł!dodatni}%
Murasugi, Przytycki zapytali na końcu pracy \cite{murasugi1989}, czy dla węzłów dodatnich rozwłóknionych (takich, których dopełnienie rozwłóknia się nad okręgiem, a~włóknem jest powierzchnia Seiferta) zachodzi równość:

\begin{conjecture}[Murasugiego-Przytyckiego]
\index[persons]{Murasugi, Kunio}%
\index[persons]{Przytycki, Józef}%
\index{hipoteza!Murasugiego-Przytyckiego}%
    Niech $K$ będzie dodatnim węzłem rozwłóknionym.
    Wtedy $\genus K = \unknotting K$.
\end{conjecture}

Dla domknięć dodatnich warkoczy równość jest znana: korzystając z~relacji w~grupie warkoczy, można rozwiązać taki węzeł ciągiem $\genus K$ zmian skrzyżowań \cite{boileau1984}.
\index[persons]{Boileau, Michel}%
\index[persons]{Weber, Claude}%
Tego rozumowania nie potrafimy uogólnić na wszystkie dodatnie węzły rozwłóknione.
W~przeciwną stronę Stojmenow zapytał, czy dla każdego alternującego węzła rozwłóknionego zachodzi $\genus K \ge \unknotting K$.
\index[persons]{Stojmenow, Aleksander}%
Według bazy KnotInfo \cite{knotinfo2024} nierówność ta jest spełniona przez wszystkie 2106 takich węzłów o~co najwyżej 13 skrzyżowaniach, przy czym jest ostra we wszystkich przypadkach poza trzynastoma \cite[problem 1.5]{baykur2026}.
